\documentclass[../DoAn.tex]{subfiles}
\begin{document}

\begin{center}
    \Large{\textbf{ABSTRACT}}\\
\end{center}
\vspace{1cm}
In public blockchain systems like Ethereum, all transactions are recorded transparently and are publicly accessible. While this characteristic ensures the transparency and verifiability of the network, it also exposes limitations regarding user privacy, as the entire transaction history, asset balances, and associations between wallet addresses can be traced. Currently, several solutions such as mixers or privacy-focused blockchains have been deployed to conceal transaction information; however, these solutions often face barriers related to integration capabilities, deployment costs, or compatibility with the Ethereum ecosystem. In this context, the Stealth Address mechanism is evaluated as a viable technical solution to enhance receiver anonymity by generating a unique receiving address for each transaction.

This project focuses on the approach of integrating Stealth Addresses with Account Abstraction on Ethereum, aiming to leverage the customizability of smart wallets while enhancing privacy in digital asset transactions. Based on this, the system is designed to allow receivers to register their public keys, senders to generate corresponding stealth addresses for each transaction, and receivers to discover and manage their assets through smart wallets.

The main contribution of this project is proposing and developing an integrated model of Stealth Addresses and Account Abstraction on the Ethereum platform, thereby supporting more private transactions while maintaining compatibility with the existing blockchain infrastructure. The achieved result is an experimental system demonstrating the applicability of the Stealth Address mechanism in protecting user privacy on the Ethereum network.

\begin{flushright}
Student\\[1cm]
\textit{(Signature and full name)}\\[2cm]

An\\
Nguyen Van An
\end{flushright}

\end{document}